%!TEX root = ../main_thesis.tex

\chapter{Introduzione}
\label{cap:introduzione}

\section{Contesto e motivazione}
\label{sec:contesto}

L'ottimizzazione matematica è uno strumento fondamentale per affrontare problemi
decisionali in ambito industriale, logistico ed economico: dalla pianificazione
della produzione all'allocazione di risorse, molte situazioni reali possono
essere ricondotte alla ricerca della soluzione migliore tra molte alternative,
nel rispetto di un insieme di vincoli. La \emph{Ricerca Operativa} e
l'\emph{Ottimizzazione Combinatoria}, da cui peraltro deriva il nome del progetto
ROOC, studiano proprio i metodi per formulare e risolvere problemi di questo tipo.

Tra la formulazione di un problema e la sua risoluzione si colloca però
un'attività non banale: la \emph{modellazione}, cioè la traduzione del problema
in un modello matematico formale che un risolutore (\emph{solver}) sia in grado
di trattare. Questa traduzione richiede competenze specialistiche e strumenti
adeguati, e gli ambienti esistenti, pur potenti, sono spesso commerciali, legati
a un linguaggio di programmazione ospite o poco orientati all'apprendimento.

A queste considerazioni se ne aggiunge una legata all'ecosistema di
\emph{Rust}~\cite{rust2019}, un linguaggio in crescita, da diversi anni il più
apprezzato dagli sviluppatori secondo il sondaggio annuale di Stack
Overflow~\cite{sosurvey2025}. Uno dei suoi punti di forza è la portabilità, dato che
lo stesso sorgente può essere compilato per molte piattaforme diverse, compreso il
browser tramite WebAssembly. Perché un'applicazione ne benefici, però, ogni sua
dipendenza deve essere \emph{pure Rust}, cioè scritta interamente in Rust e
priva di dipendenze native. Gli strumenti di modellazione più diffusi non
soddisfano questo requisito, perché richiedono un runtime esterno, come
l'interprete Python per Pyomo o l'eseguibile di AMPL, oppure librerie native
precompilate, come nel caso di Gurobi. D'altra parte, al momento della
scrittura nessuno degli strumenti di modellazione nati nell'ecosistema Rust
offre vincoli logici su modelli misti né la loro linearizzazione automatica
(Sezione~\ref{sec:stato-arte}).

\section{Il problema e l'idea di ROOC}
\label{sec:idea}

ROOC nasce per rendere più accessibile l'attività di
modellazione e colmare questa lacuna dell'ecosistema
Rust.\footnote{Il codice sorgente del progetto è disponibile
pubblicamente all'indirizzo \url{https://github.com/Specy/rooc}.} È un
linguaggio dichiarativo, con una notazione vicina a quella matematica, che
consente di descrivere un modello di ottimizzazione separandone la struttura dai
dati. A partire da tale descrizione, ROOC si comporta come un \emph{compilatore}:
attraverso una sequenza di trasformazioni traduce il modello di alto livello in
una forma lineare standard, che viene poi affidata a un solver per ottenere la
soluzione ottima.

L'idea di fondo è dunque applicare le tecniche dei compilatori al dominio
dell'ottimizzazione, ottenendo uno strumento che non solo risolve i modelli, ma
ne rende ispezionabile ogni fase intermedia. Questa trasparenza, unita
all'esecuzione interamente nel browser, conferisce a ROOC anche una valenza
didattica.

Il motore è inoltre pensato per essere incorporato in altre applicazioni. È
\emph{pure Rust}, solver predefiniti compresi, e oltre al linguaggio testuale
offre una \emph{fluent API}, con cui un modello si costruisce direttamente nel
codice Rust. Può quindi essere compilato per qualunque target supportato dal
compilatore Rust, e la piattaforma web, eseguita interamente nel browser, ne è
una dimostrazione.

\section{Obiettivi e contributi}
\label{sec:contributi}

Il presente lavoro si concentra sul motore di ROOC, e in particolare sulla
progettazione del linguaggio e del relativo compilatore. I contributi principali
sono:

\begin{itemize}
  \item la progettazione di un \emph{linguaggio} di modellazione dotato di un
        sistema di tipi e di costrutti di alto livello, quali iteratori,
        \emph{block function} e tipi di dato strutturati
        (Capitolo~\ref{cap:linguaggio});
  \item la realizzazione di un \emph{compilatore multi-fase}, che traduce il
        modello fino alla forma standard, organizzato secondo un'architettura a
        \emph{pipe} che rende ogni fase componibile e ispezionabile
        (Capitoli~\ref{cap:architettura}--\ref{cap:linearizzazione});
  \item l'estensione del solver \texttt{microlp}, ottenuto come \emph{fork} del
        crate archiviato \texttt{minilp} e ampliato dal solo simplesso al
        supporto delle variabili intere e binarie, fino a coprire la
        programmazione lineare intera mista (Capitolo~\ref{cap:solver});
  \item la distribuzione del motore, pure Rust insieme ai suoi solver
        predefiniti, come libreria Rust, che accanto al
        linguaggio offre una \emph{fluent API} per costruire i modelli
        direttamente nel codice, come pacchetto TypeScript precompilato e come
        piattaforma web a fini didattici (Capitolo~\ref{cap:piattaforma}).
\end{itemize}
